This proposal extends the e-graph with semantic information beyond the
equalities it was designed to represent, into a structure that reasons about a
broader class of relationships between terms. Taking \textit{proof by cases} as
our guide, we make disequalities and conditional facts first-class citizens of
the e-graph and exploit the resulting structure for optimization. Since the
e-graph is the shared substrate of much of the theorem-proving ecosystem, these
extensions stand to benefit automated and interactive provers alike.

Figure~\ref{fig:plan} outlines the high-level plan for the five years of the
Ph.D., split into two phases. During the \textit{course phase}, in the first two
years, we develop the two \textit{semantic} extensions that enrich what the
e-graph represents: disequality reasoning and conditional reasoning. During the
\textit{dissertation phase}, we pursue the \textit{optimizations} that these
semantics make possible: lazy reasoning and parallel reasoning. During the
development, the two lines converge into one final deliverable,
\textit{Smart E-Graphs}, which we will integrate into both automated
and interactive provers to demonstrate the benefits of our work.

\begin{figure}[t]
\centering
\definecolor{darkgreen}{HTML}{2A5E32}
\definecolor{midgreen}{HTML}{377E39}

\begin{tikzpicture}[
  yearline/.style={gray, opacity=0.35, thin, dashed},
  chevron/.style={
    signal, signal to=east, signal from=west, signal pointer angle=110,
    draw=black, line width=0.6pt, inner sep=6pt, text=white,
    align=center, minimum height=1.3cm
  },
  outer/.style={chevron, fill=darkgreen, minimum height=3.1cm},
  inner/.style={chevron, fill=midgreen, minimum width=3.2cm},
  label/.style={font=\bfseries}
]

% ---- Time Axis ----
\foreach \x/\lab in {-3.25/0y, -0.01/1y, 3.23/2y, 6.47/3y, 9.71/4y, 12.95/5y}{
  \draw[yearline] (\x,-2.3) -- (\x,2.0);
  \node[gray, font=\small\bfseries] at (\x,2.25) {\lab};
}

% ---- Course Phase ----
\node[outer, minimum width=0.38\linewidth] (g1) at (0,0) {};
\node[inner, text width=0.15\linewidth] at ($(g1.center)+(0,0.75)$) {Disequality\\Reasoning};
\node[inner, text width=0.15\linewidth] at ($(g1.center)+(0,-0.75)$) {Conditional\\Reasoning};
\node[label] (g1label) at ($(g1.center)+(0,-2.8)$) {Course Phase};

% ---- Dissertation Phase ----
\node[outer, minimum width=0.57\linewidth] (g2) at (8.1,0) {};
\node[inner, text width=0.15\linewidth] at ($(g2.center)+(-1.6,0.75)$) {Lazy\\Reasoning};
\node[inner, text width=0.15\linewidth] at ($(g2.center)+(-1.6,-0.75)$) {Parallel\\Reasoning};
\node[chevron, fill=midgreen, text width=0.10\linewidth, minimum height=2.8cm]
     (tp) at ($(g2.center)+(2.3,0)$) {};
\node[align=center, text=white] at ([xshift=2mm]tp.center) {Smart\\E-Graphs};
\node[label] (g2label) at ($(g2.center)+(0,-2.8)$) {Dissertation Phase};
\end{tikzpicture}
\caption{High-level plan for the Ph.D. proposal.}
\label{fig:plan}
\end{figure}